\documentclass[11pt]{amsart}
\usepackage[T1]{fontenc}
\usepackage{lmodern}
\usepackage{amsmath,amssymb,amsthm}
\usepackage[numbers,sort]{natbib}
\usepackage[hidelinks]{hyperref}
\usepackage{bm}

\newtheorem{theorem}{Theorem}
\newtheorem{lemma}[theorem]{Lemma}
\newtheorem{corollary}[theorem]{Corollary}
\theoremstyle{definition}
\newtheorem{definition}[theorem]{Definition}
\newtheorem{remark}[theorem]{Remark}

\newcommand{\Ter}{\mathcal S}
\newcommand{\app}{\mathbin{+\!\!+}}
\newcommand{\args}{\mathrm{args}}
\newcommand{\tw}[1]{\mathsf{tw}_{#1}}
\newcommand{\run}{\mathrm{run}}
\newcommand{\seat}{\mathrm{seat}}
\newcommand{\Fires}{\mathrm{Fires}}
\newcommand{\Silent}{\mathrm{Silent}}
\newcommand{\Eternal}{\mathrm{Eternal}}
\newcommand{\Quiet}{\mathrm{Quiet}}
\newcommand{\QuietD}{\mathrm{Quiet}^{\exists}}
\newcommand{\BOT}{\mathrm{BOT}}
\newcommand{\KCAP}{\mathrm{KCAP}}
\newcommand{\RD}{\mathrm{RD}}
\newcommand{\FOD}{\mathrm{FOD}}
\newcommand{\capf}{\mathrm{cap}}
\newcommand{\prog}{\mathrm{prog}}
\newcommand{\lam}{\Lambda}
\newcommand{\rd}{\mathrm{reader}}
\newcommand{\cps}{\mathrm{capsule}}
\newcommand{\pads}{\mathrm{pads}}
\newcommand{\lean}[1]{\texttt{#1}}
\newcommand{\Seq}{\mathsf{Seq}}
\newcommand{\Iter}{\mathsf{Iter}}
\newcommand{\Loop}{\mathsf{Loop}}
\newcommand{\Cond}{\mathsf{Cond}}
\newcommand{\Reindex}{\mathsf{Reindex}}
\newcommand{\Par}{\mathsf{Par}}

\title[Clock Closure]{Clock Closure on Head Runs of the Combinator $\bm{S}$, and a Uniform Countdown}
\author{Claude Opus \& Fable with N.\ E.\ Davis}
\address{ZeroAttest, Champaign, Illinois, USA}
\email{NDavis@ZeroAttest.com}
\urladdr{https://orcid.org/0000-0002-1131-8319}
\date{\today}

\begin{document}
\begin{abstract}
A head step on a closed term built from the combinator $S$ pops the first three arguments
$x,y,z$ of the head and pushes the arguments of $x$, then $z$ and $y\,z$. The step \emph{fires} if
$x$ is the normal term $F=S\,\tw8\,\tw8$, where $\tw8$ is the right comb of eight $S$'s. A
\emph{clock} is a family of argument lists together with its schedule, given as data: whether the
head run fires and after how many steps; when it does not fire, its run must also never halt. We
prove that clocks are closed under five operations on head runs: a quiet prefix followed by a
clock, iteration with an exit at the bottom, an endless loop, sums and reindexing. With these we
prove that a countdown term with a program of $m$ layers, whose tags are arbitrary terms, and a
register holding $g$ first fires at head step $22+(5m+13)g$, for every program and every $g$; two
capsules of four and six leaves first fire at steps $5g+1$ and $7g+1$. Everything is mechanized in
Lean~4, and the only axioms are propositional extensionality and quotient soundness.
\end{abstract}

\keywords{combinatory logic, combinator $S$, head reduction, clocks, Lean}
\maketitle

\section{Introduction}

$\Ter$ is the set of closed terms built from $S$ by application. Every $t\in\Ter$ is
$S\,u_1\cdots u_m$ for a unique $m\ge0$, and $\args(t)=[u_1,\dots,u_m]$ is the list of its
items~\citep{s1,s11}. If $m\ge3$, the \emph{head step} contracts the leftmost redex,
$t\to_h u_1\,u_3\,(u_2\,u_3)\,u_4\cdots u_m$, and then
\begin{equation}\label{eq:head}
\args(u_1\,u_3\,(u_2\,u_3)\,u_4\cdots u_m)=\args(u_1)\app[u_3,\;u_2\,u_3,\;u_4,\dots,u_m]
\end{equation}
\citep[equation~(1)]{s1}. So the head run of $t$ is a run on the list $\args(t)$ alone. Write $\tw n$ for the
right comb with $n$ leaves ($\tw0=\tw1=S$, $\tw{n+2}=S\,\tw{n+1}$), and let $F=S\,\tw8\,\tw8$, the
\emph{marker}, a normal form with seventeen leaves. A head step \emph{fires} if its $u_1$ is $F$,
that is, if the term is $S\,F\,y\,z\cdots$. The question is when a term first fires, as a function
of parameters built into it.

Write $SS$ for $S\,S$, and let
\begin{gather*}
A=SS\,S,\qquad \FOD=SS\,SS,\qquad g_1=S\,(S\,S\,\FOD),\qquad \RD=S\,g_1,\\
\KCAP=S\,SS\,(S\,(S\,SS)),\qquad \BOT=S\,F\,F.
\end{gather*}
A \emph{layer} is $\lam(W,s)=S\,SS\,(SS\,W)\,s$. The \emph{program} with tags $W_1,\dots,W_m$ is
$\prog(\vec W)=\lam(W_1,\prog(W_2,\dots,W_m))$, with $\prog()=A$. The \emph{register} is
$\rho_0=\BOT$, $\rho_{c+1}=\KCAP\,\rho_c$. With $\sigma=\prog(\vec W)$, the \emph{countdown
reader} is
\[
\rd(\vec W,g)=S\,S\,(\rho_g\,(\RD\,\sigma))\,\tw7\,(S\,\tw4\,\tw7)\,\tw8 .
\]

\begin{theorem}\label{thm:main}
Let $W_1,\dots,W_m\in\Ter$, $g\ge0$ and $N=21+(5m+13)g$. The head run of $\rd(\vec W,g)$ makes at
least $N$ steps; its term after $N$ steps is $S\,F\,y\,z\cdots$, and no earlier term of the run is.
So the first firing is head step $22+(5m+13)g$.
\end{theorem}

We count steps from $1$: the step that fires is step $N+1$, and its \emph{firing index} is $N+1$.
The Lean statements count the $N$ steps before it (0-based). Every index below is quoted 1-based.

The proof is one instance of a small algebra of head runs. Section~\ref{sec:frames} sets up runs
on argument lists with an arbitrary tail; Section~\ref{sec:closure} defines clocks and proves their
closure under five operations; Section~\ref{sec:capsules} proves two countdowns of one revolution
each by hand; Section~\ref{sec:reader} proves Theorem~\ref{thm:main}. The steps of a run are
computed by the Lean kernel with the tags, the program and the rest of the register left as free
variables, so the theorem holds for every program at once, not word by word. Section~\ref{sec:scope}
says what is not proved.

\section{Frames and firing}\label{sec:frames}

A \emph{frame} is a finite list of terms. The \emph{step} on a frame is
\[
[x,y,z]\app r\;\mapsto\;\args(x)\app[z,\;y\,z]\app r,
\]
undefined on frames with fewer than three entries; then $x$ is the \emph{seat}. Let $\run_n(l)$ be
the frame after $n$ steps from $l$, if they are all defined. By~\eqref{eq:head}, the head run of
$t$ is the run of $\args(t)$, and the frame of $t$ is $\args(t)$.

\begin{definition}\label{def:fire}
Let $l,a,b$ be frames and $n\ge0$.
\begin{itemize}
\item $\Fires_n(l)$: $\run_n(l)$ is defined and has seat $F$, and for no $k<n$ is
  $\run_k(l)$ defined with seat $F$.
\item $\Silent(l)$: no defined $\run_k(l)$ has seat $F$.
\item $\Eternal(l)$: every $\run_k(l)$ is defined and has at least three entries.
\item $\Quiet_n(a,b)$: for every frame $R$, $\run_n(a\app R)=b\app R$ and no $\run_k(a\app R)$
  with $k<n$ has seat $F$. $\QuietD_n(a,b)$: $\Quiet_n(a,b\app J)$ for some $J$, the
  \emph{deposit}.
\end{itemize}
\end{definition}

A quiet segment is stated for every tail $R$, and its deposit stays behind the window of three
entries. $\Fires_n(l)$ says that step $n+1$ is the first firing.

\begin{lemma}\label{lem:frame}
\textup{(a)} If $\run_n(l)=l'$, then $\run_n(l\app R)=l'\app R$.

\textup{(b)} Let $\run_n(a)=b$, and let no $\run_k(a)$ with $k<n$ have seat $F$. Then
$\Quiet_n(a,b)$.

\textup{(c)} If $\Quiet_n(a,b\app J)$ and $\Quiet_m(b,c)$, then $\Quiet_{n+m}(a,c\app J)$.

\textup{(d)} If $\Quiet_n(a,b)$ and $\Fires_m(b\app R)$, then $\Fires_{n+m}(a\app R)$.

\textup{(e)} If $\Fires_n(l)$, then $\Fires_n(l\app R)$ and not $\Silent(l)$; if also
$\Fires_m(l)$, then $n=m$.

\textup{(f)} If $\Quiet_n(a,b)$ and $b$ is silent and eternal, so is $a$. If $l$ is silent and
eternal, so is $l\app J$.
\end{lemma}

\begin{proof}
(a) A defined step reads only the first three entries; induction on $n$. (b) By (a), the frames
from $a\app R$ are those from $a$ with $R$ appended, and each frame before the last has three
entries, since it steps, so its seat is unchanged. (c) Apply $\Quiet_m(b,c)$ to the tail $J\app R$.
(d) Split each $k<n+m$ at $n$. (e) By (b), $a=l$ runs quietly to $\run_n(l)$, whose seat survives
appending; (d) gives the first claim, the others follow from minimality. (f) The first claim splits
$k$ at $n$. Every frame on an eternal run has three entries, so by (a) appending $J$ changes no
seat and no length below three.
\end{proof}

Lemma~\ref{lem:frame} holds for lists of arbitrary terms. To read a frame run as head reduction we
need the head symbol $S$, which $\args$ drops.

\begin{lemma}\label{lem:bridge}
Let $t\in\Ter$ and $n\ge0$. The head run of $t$ makes at least $n$ steps iff $\run_n(\args(t))$ is
defined, and then its term after $n$ steps is $S\,v_1\cdots v_k$ with
$[v_1,\dots,v_k]=\run_n(\args(t))$. Hence, if $\Fires_n(\args(t))$, the term after $n$ head
steps is $S\,F\,y\,z\cdots$ and no earlier term of the run is.
\end{lemma}

\begin{proof}
By~\eqref{eq:head} and induction on $n$; the steps preserve membership in $\Ter$, and every
$u\in\Ter$ is $S$ applied to $\args(u)$.
\end{proof}

\emph{Checked runs.} Let $\run^F_n$ be the run that is undefined as soon as a seat is $F$. If
$\run^F_n(a)=b$, then $\Quiet_n(a,b)$, by Lemma~\ref{lem:frame}(b); if moreover $b$ has seat $F$,
then $\Fires_n(a\app R)$ for every $R$. The test ``seat $=F$'' is a pattern match on
$S\,\tw8\,\tw8$, and the step needs $\args(x)$ only for the seat $x$. So the Lean kernel can
evaluate $\run^F_n(a)$ when $a$ contains free variables, as long as no free variable is, or heads,
a seat; if one were, the evaluation would be stuck and the proof would fail. Every segment lemma
below is one such evaluation, and its free variables are the terms it quantifies over.

\section{Clocks and their closure}\label{sec:closure}

\begin{definition}\label{def:clock}
Let $P$ be a type. A \emph{clock} on $P$ consists of frames $\phi(p)$, a predicate $H\colon P\to
\{\mathrm{true},\mathrm{false}\}$ and a function $T\colon P\to\mathbb N$, with proofs that
\begin{itemize}
\item if $H(p)$, then $\Fires_{T(p)}(\phi(p)\app R)$ for every $R$;
\item if not $H(p)$, then $\phi(p)$ is silent and eternal.
\end{itemize}
\end{definition}

The pair $(T,H)$ is the \emph{schedule}. The quiet side asks for eternity as well as silence. This
is stricter than the dynamical definition of~\citep{clocknote}: a family whose run halts without
firing, such as the empty frame, is not a clock. Eternity is what lets a quiet side pass through a
deposit (Lemma~\ref{lem:frame}(f)). The tail $R$ in the firing side follows from
Lemma~\ref{lem:frame}(e) and is kept for convenience.

\begin{theorem}\label{thm:closure}
\textup{($\Seq$)} Let $C$ be a clock on $Q$ with frames $\phi$. For $p\in P$ let $a(p)$ and
$J(p)$ be frames, $o(p)\in Q$ and $c(p)\in\mathbb N$, and let them satisfy $\Quiet_{c(p)}(a(p),\phi(o(p))\app J(p))$ for
all $p$. Then for every $p$: if $H(o(p))$, then $\Fires_{c(p)+T(o(p))}(a(p)\app R)$ for every $R$;
otherwise $a(p)$ is silent and eternal. So $a$ is a clock on $P$ with $H\circ o$ and
$c+T\circ o$.

\textup{($\Iter$)} For $c\in\mathbb N$ let $b(c)$ and $J(c)$ be frames and $\pi(c)\in\mathbb N$;
let $e\in\mathbb N$, and let them satisfy $\Quiet_{\pi(c)}(b(c+1),b(c)\app J(c))$ for all $c$, and
$\Fires_e(b(0)\app R)$ for all $R$. Then $\Fires_{e+\sum_{c<g}\pi(c)}(b(g)\app R)$ for all $g$
and $R$. So $b$ is a clock on $\mathbb N$ with $H\equiv\mathrm{true}$ and
$T(g)=e+\sum_{c<g}\pi(c)$.

\textup{($\Loop$)} For $k\in\mathbb N$ let $\kappa(k)$ and $J(k)$ be frames and
$\pi(k)\in\mathbb N$, and let them satisfy $\Quiet_{\pi(k)+1}(\kappa(k),\kappa(k+1)\app J(k))$ for all $k$. Then $\kappa(0)$ is silent
and eternal.

\textup{($\Cond$)} For clocks $C_1$ on $P_1$ and $C_2$ on $P_2$ there is a clock on $P_1\sqcup P_2$
whose frames, $T$ and $H$ are those of $C_1$ on $P_1$ and those of $C_2$ on $P_2$.

\textup{($\Reindex$)} For a clock $C$ on $Q$ with frames $\phi$ and a function $f\colon P\to Q$
there is a clock on $P$ with frames $\phi\circ f$, $T\circ f$ and $H\circ f$.
\end{theorem}

\begin{proof}
($\Seq$) Apply Lemma~\ref{lem:frame}(d) with the tail $J(p)\app R$, and on the quiet side
both parts of Lemma~\ref{lem:frame}(f). ($\Iter$) Induction on $g$: by Lemma~\ref{lem:frame}(d)
with the tail $J(g)\app R$, the segment from $b(g+1)$ adds $\pi(g)$ to the firing of $b(g)$; this
is why the exit is assumed for every tail. ($\Loop$) By strong induction on $j$, for all $k$ and
$R$ the frame $\run_j(\kappa(k)\app R)$ is defined, has three entries and a seat other than $F$:
if $j\le\pi(k)$ the segment from $\kappa(k)$ says so; otherwise its run reaches
$\kappa(k+1)\app J(k)\app R$ after $\pi(k)+1\ge1$ steps, and the induction hypothesis applies with
$k+1$ and $j-\pi(k)-1$. Take $k=0$ and $R$ empty. ($\Cond$) By cases on the sum. ($\Reindex$) By
composition.
\end{proof}

A conditional ``if $\theta(p)$ then $C_1$ else $C_2$'' with $\theta$ a function to Booleans is
$\Reindex$ of the sum along $p\mapsto$ ($\iota_1 p$ if $\theta(p)$, else $\iota_2 p$). A loop needs
a family $\kappa$, not one frame: if $\kappa(k)=l$ for all $k$ and $l$ is a frame of $S$-terms,
then by Lemma~\ref{lem:bridge} the term $S\,l_1\cdots l_r$ would head-reduce in at least one step
to itself applied to the deposit, which no closed $S$-term does~\citep[Corollary~2]{s7}.

\begin{corollary}\label{cor:iff}
For a clock and every $p$: $\Fires_n(\phi(p))$ for some $n$ iff $H(p)$, and then $n=T(p)$.
\end{corollary}

\begin{proof}
If not $H(p)$, $\phi(p)$ is silent, and Lemma~\ref{lem:frame}(e) excludes a firing. If $H(p)$, take
$R$ empty; uniqueness is Lemma~\ref{lem:frame}(e).
\end{proof}

\begin{remark}\label{rem:scope}
A clock is a schedule paired with a proof. Nothing in Definition~\ref{def:clock} makes $T$ or $H$
computable, and Corollary~\ref{cor:iff} says only that $H$ is the firing predicate and $T$ the
first firing. Computability of the schedule, and so decidability of firing, holds only for the
named instances \lean{k4Clock}, \lean{c6Clock} and \lean{readerClock} of
Sections~\ref{sec:capsules} and~\ref{sec:reader}, whose $T$ is the closed form $5g$, $7g$ or
$21+(5m+13)g$ and whose $H$ is constantly true. The quiet side ($\Loop$
and the quiet halves of $\Seq$ and $\Cond$) is proved, but no clock in this paper uses it: no
family that fires only for some parameters, such as a gate that fires iff $g=0$, is proved here.
\end{remark}

\section{Two capsule countdowns}\label{sec:capsules}

Let $K_4=S\,A=S\,(S\,S\,S)$ and $C_6=S\,(S\,(S\,A))$. For a term $C$ let $C^{(0)}=\BOT$,
$C^{(g+1)}=C\,C^{(g)}$, and
\begin{gather*}
\cps(C,g)=C^{(g)}\,A\,\tw7\,(S\,\tw4\,\tw7)\,\tw8\,\tw5,\\
\pads=[\tw7,\;S\,\tw4\,\tw7,\;\tw8,\;\tw5],\\
\args(\cps(C,0))=[F,F,A]\app\pads,\\
\args(\cps(K_4,g+1))=[A,K_4^{(g)},A]\app\pads.
\end{gather*}
The first of these frames fires at once.

\begin{lemma}\label{lem:k4}
For every term $V$ we have $\QuietD_5([A,K_4\,V,A],[A,V,A])$, and
$\Fires_5([A,\BOT,A]\app R)$ for every $R$.
\end{lemma}

\begin{proof}
Let $X$ be $K_4\,V$ or $\BOT$. As $\args(A)=[S,S]$ and $S\,A=K_4$,
\begin{align*}
[A,X,A]&\to[S,S,A,X\,A]\to[A,K_4,X\,A]\\
&\to[S,S,X\,A,K_4\,(X\,A)]\to[X\,A,\;S\,(X\,A),\;K_4\,(X\,A)].
\end{align*}
The seats are $A$, $S$, $A$, $S$ and $X\,A$. None is $F$: the items of $A$ are $S,S$, and $X\,A$
has three items. If $X=K_4\,V$, then $\args(X\,A)=[A,V,A]$, and the fifth step gives $[A,V,A]$
followed by a deposit. If $X=\BOT$, then $\args(X\,A)=[F,F,A]$: the fifth step leads to
a frame with seat $F$.
\end{proof}

The run never opens $V$, which is why $V$, the rest of the register, may be any term.

\begin{theorem}\label{thm:capsules}
For every $g\ge0$,
\[
\Fires_{5g}\bigl(\args(\cps(K_4,g))\bigr)\quad\text{and}\quad\Fires_{7g}\bigl(\args(\cps(C_6,g))\bigr).
\]
So the head run of $\cps(K_4,g)$ first fires at step $5g+1$, and that of $\cps(C_6,g)$ at step $7g+1$.
\end{theorem}

\begin{proof}
For $g=0$ the frame is $[F,F,A]\app\pads$. Otherwise apply $\Iter$ to $b(c)=[A,K_4^{(c)},A]$,
$\pi\equiv5$, $e=5$ (Lemma~\ref{lem:k4}): the frame $b(g-1)\app\pads$ fires after
$5+5(g-1)=5g$ steps. For $C_6$ the boundary is $[S\,K_4,C_6^{(c)},A]$, and one kernel evaluation
gives $\QuietD_7([S\,K_4,C_6\,V,A],[S\,K_4,V,A])$ for every $V$ and
$\Fires_7([S\,K_4,\BOT,A]\app R)$ for every $R$.
Lemma~\ref{lem:bridge} gives the indices, since $\cps(C,g)\in\Ter$ for $C\in\{K_4,C_6\}$.
\end{proof}

\section{The countdown reader}\label{sec:reader}

We use the terms of the introduction. Write $\sigma=\prog(\vec W)$ and
$\capf(\sigma,c)=[g_1,\;\sigma,\;\rho_c\,(\RD\,\sigma)]$, the \emph{capture} frame. The tag
terms $W_i$ are arbitrary.

\begin{lemma}[layer]\label{lem:layer}
For all terms $W,W',s,Q$, with $\lam'=\lam(W',s)$,
\[
\Quiet_5\bigl([S,W,\lam',Q],\;[S,W',s,Q]\app[\,s\,(SS\,W'\,s)\,Q,\;W\,\lam'\,Q\,]\bigr).
\]
\end{lemma}

\begin{proof}
As $\args(\lam')=[SS,\,SS\,W',\,s]$, $\args(SS)=[S]$ and $\args(SS\,W'\,s)=[S,W',s]$,
\begin{align*}
[S,W,\lam',Q]&\to[\lam',\;W\lam',\;Q]\to[SS,\;SS\,W',\;s,\;Q,\;W\lam'\,Q]\\
&\to[S,\;s,\;SS\,W'\,s,\;Q,\;W\lam'\,Q]\\
&\to[SS\,W'\,s,\;s\,(SS\,W'\,s),\;Q,\;W\lam'\,Q]\\
&\to[S,\;W',\;s,\;Q,\;s\,(SS\,W'\,s)\,Q,\;W\lam'\,Q].
\end{align*}
The seats $S$, $\lam'$, $SS$, $S$, $SS\,W'\,s$ have $0$, $3$, $1$, $0$ and $3$ items, and $F$
has $2$, so none is $F$. The tail behind the window is untouched (Lemma~\ref{lem:frame}(b)).
\end{proof}

No free variable is a seat, so $W,W',s,Q$ are arbitrary. The step with seat $S$ and the tag $W$
in the next entry opens each layer.

\begin{lemma}[segments]\label{lem:segments}
Let $W,s,Q,C,\sigma$ be terms, $R$ a frame, and
$t(X)=S\,S\,(X\,(\RD\,\sigma))\,\tw7\,(S\,\tw4\,\tw7)\,\tw8$. Then:
\begin{enumerate}
\item[(a)] $\QuietD_{28}\bigl(\args(t(\KCAP\,C)),\;[g_1,\sigma,C\,(\RD\,\sigma)]\bigr)$;
\item[(b)] $\Fires_{21}\bigl(\args(t(\BOT))\app R\bigr)$;
\item[(c)] $\QuietD_7\bigl([g_1,\lam(W,s),Q],\;[S,W,s,Q]\bigr)$;
\item[(d)] $\QuietD_{11}\bigl([S,W,A,\KCAP\,C\,(\RD\,\sigma)],\;
  [g_1,\sigma,C\,(\RD\,\sigma)]\bigr)$;
\item[(e)] $\Fires_4\bigl([S,W,A,\BOT\,(\RD\,\sigma)]\app R\bigr)$;
\item[(f)] $\QuietD_{13}\bigl([g_1,A,\KCAP\,C\,(\RD\,A)],\;[g_1,A,C\,(\RD\,A)]\bigr)$,
  and\\ $\Fires_6\bigl([g_1,A,\BOT\,(\RD\,A)]\app R\bigr)$.
\end{enumerate}
\end{lemma}

\begin{proof}
Each is one kernel evaluation of a checked run (Section~\ref{sec:frames}) with the displayed
terms free. In (c) the run opens the layer $\lam(W,s)$ but not $W$ or $s$; in (d) and (f) it
opens $\KCAP\,C$ but not $C$, and it opens $\RD\,\sigma$ but not $\sigma$.
\end{proof}

The segment (c) starts a revolution: it opens the first layer of the program. The segment (d)
ends one: from the base $A$ it pops the top $\KCAP$ of the register and rebuilds the capture frame
with the register decremented. Part (f) is the case $m=0$, where the program is $A$ and (c) does
not apply.

\begin{lemma}\label{lem:rev}
Let $m=|\vec W|$ and $\sigma=\prog(\vec W)$. For every $c$ and every $R$,
\[
\QuietD_{5m+13}\bigl(\capf(\sigma,c+1),\capf(\sigma,c)\bigr)\quad\text{and}\quad
\Fires_{5m+6}\bigl(\capf(\sigma,0)\app R\bigr).
\]
\end{lemma}

\begin{proof}
By Lemma~\ref{lem:layer}, Lemma~\ref{lem:frame}(c) and induction on the word,
\[
\QuietD_{5k}\bigl([S,W,\prog(\vec V),Q],\;[S,W',A,Q]\bigr)
\]
for every word $\vec V$ of length $k$, where $W'$ is the last letter of $W\vec V$. Let $m\ge1$.
Chain Lemma~\ref{lem:segments}(c) for the first layer, these $5(m-1)$ steps for the others, and
then (d) with $Q=\rho_{c+1}\,(\RD\,\sigma)$, or (e) with $Q=\BOT\,(\RD\,\sigma)$. This gives $7+5(m-1)+11=5m+13$ and
$7+5(m-1)+4=5m+6$ (Lemma~\ref{lem:frame}(c),(d)). For $m=0$ use (f): $13$ and $6$.
\end{proof}

\begin{proof}[Proof of Theorem~\ref{thm:main}]
By $\Iter$ with $b(c)=\capf(\sigma,c)$, $\pi\equiv5m+13$ and $e=5m+6$ (Lemma~\ref{lem:rev}),
\[
\Fires_{(5m+6)+(5m+13)c}\bigl(\capf(\sigma,c)\app R\bigr)\qquad\text{for all $c$ and $R$.}
\]
For $g=0$, Lemma~\ref{lem:segments}(b) gives $\Fires_{21}(\args(\rd(\vec W,0)))$. For $g=c+1$,
$\rho_g=\KCAP\,\rho_c$, so Lemma~\ref{lem:segments}(a) with $C=\rho_c$ runs
$\args(\rd(\vec W,g))$ in $28$ quiet steps to $\capf(\sigma,c)$ and a deposit, and
Lemma~\ref{lem:frame}(d) gives a first firing after
\[
28+(5m+6)+(5m+13)c=21+(5m+13)(c+1)=21+(5m+13)g
\]
steps. As $\rd(\vec W,g)\in\Ter$, Lemma~\ref{lem:bridge} turns this into the statement on
head reduction.
\end{proof}

The proof gives $\Fires_{21+(5m+13)g}(\args(\rd(\vec W,g))\app R)$ for every tail $R$ and, in the
mechanization, also for tags that contain $K$ or $I$. So the reader is a clock on pairs
$(\vec W,g)$ with $H$ constantly true and $T=21+(5m+13)g$.

\begin{remark}\label{rem:reads}
On its way to the firing the run passes frames that begin with $\capf(\sigma,c)$, for
$c=g-1,\dots,0$, and in each revolution frames that begin with $[S,W_i]$, for $i=1,\dots,m$ in
order. This is a feature of
the proof. We do not prove that the run passes exactly $g$ capture frames, or that the steps with
seat $S$ and a tag in the next entry spell $W_1\cdots W_m$ exactly $g$ times. So ``the countdown
reads its program'' is a description, not a theorem.
\end{remark}

\section{Scope}\label{sec:scope}

The operations are those of Theorem~\ref{thm:closure}: $\Seq$ is a quiet prefix followed by a
clock; $\Iter$ has a single exit at the bottom; $\Cond$ is a sum; $\Reindex$ pulls a clock back
along a function; and $\Loop$ is new. Closure is proved operation by operation. There is no
theorem about a least class closed under them, and the statement of~\citep[Theorem~2]{clocknote},
that no composite witnesses universality, is not mechanized and is not claimed here. A product
$\Par$, running two machines ``in disjoint frame regions'', is dropped: a head run moves one window
over one frame, so there are no two regions that evolve at once, and an interleaving of two
machines within one head run is a $\Seq$ or an $\Iter$.

The other schedules of the clock shop~\citep{murm2}, namely the transfer train ($6n+3$), the fold
apparatus ($26n+5$), the ratio encoder ($26qn+4$ iff $g=0$), the two-counter relay
($12a+6b+4$ iff $g=0$) and the gate, are observed on sampled parameters. They are not proved, here
or there.

\emph{Corrections to earlier claims.} The row ``countdown reader, always, $21+23g$''
of~\citep[Proposition~1]{clocknote} is right only for $m=2$, and only as a 0-based count; by
Theorem~\ref{thm:main} the first firing is step $22+(5m+13)g$, as already observed
in~\citep{murm2}. The schedules of that proposition were called verified over their classes; they
were sampled, not proved. For the reader, \citep{murm2} gives a symbolic execution, checked in Python, for each of the
fifteen words of length at most three over two tag constants; Theorem~\ref{thm:main} holds for
every program.

\section*{Data availability}
\sloppy
%% TODO at release: replace "branch s10-clock, commit df48181 (to be released as v1.22.0)"
%% by the release version and commit (and update refs.bib).
Every numbered result above is mechanized in \texttt{lean-ski}~\citep{leanski} (Lean~4~\citep{lean4}),
branch \texttt{s10-clock}, commit \texttt{df48181} (to be released as v1.22.0), in the modules
\texttt{Ski/S10Frame.lean}, \texttt{S10Closure.lean}, \texttt{S10Capsule.lean} and
\texttt{S10Reader.lean}, and uses only the axioms \texttt{propext} and \texttt{Quot.sound}; the gate
\texttt{Ski/S10Audit.lean} checks every declaration of these modules. Lean's terms also contain
$K$ and $I$; the frame-level statements hold for arbitrary terms, and the term-level ones assume
\texttt{PureS}. Steps are counted 0-based.
Definition~\ref{def:fire} is \texttt{S10.fstep}, \texttt{S10.frun}, \texttt{S10.seat},
\texttt{S10.FiresAt}, \texttt{S10.Silent}, \texttt{S10.Eternal}, \texttt{S10.Quiet} and
\texttt{S10.QuietD}, with \texttt{S10.args}, \texttt{S10.tw} and \texttt{S10.F}.
Lemma~\ref{lem:frame} is \texttt{S10.frun\_append}, \texttt{S10.Quiet.of\_nil},
\texttt{S10.Quiet.trans}, \texttt{S10.Quiet.fires}, \texttt{S10.FiresAt.append},
\texttt{S10.FiresAt.not\_silent}, \texttt{S10.FiresAt.unique},
\texttt{S10.Quiet.silent\_eternal} and \texttt{S10.silent\_eternal\_append}.
Lemma~\ref{lem:bridge} is \texttt{S10.iter\_eq\_frun} and \texttt{S10.TFiresAt.of\_firesAt}, and
checked runs are \texttt{S10.qrun}, \texttt{S10.Quiet.of\_qrun}, \texttt{S10.Quiet.of\_qrun\_take}
and \texttt{S10.firesAt\_of\_qrun}. Definition~\ref{def:clock} is \texttt{S10.Clock};
Theorem~\ref{thm:closure} is \texttt{S10.clock\_closure}, with the constructors
\texttt{S10.Clock.seq}, \texttt{S10.Clock.iter}, \texttt{S10.Clock.cond} and
\texttt{S10.Clock.reindex} and the lemmas \texttt{S10.iter\_fires} and
\texttt{S10.loop\_silent\_eternal}; Corollary~\ref{cor:iff} is \texttt{S10.Clock.fires\_iff} and
\texttt{S10.Clock.firesAt\_T}. Lemma~\ref{lem:k4} is \texttt{S10.k4\_rev} and
\texttt{S10.k4\_exit}; Theorem~\ref{thm:capsules} is \texttt{S10.k4\_fires}, \texttt{S10.c6\_fires}
(with \texttt{S10.c6\_rev}, \texttt{S10.c6\_exit}), \texttt{S10.k4\_fires\_term},
\texttt{S10.c6\_fires\_term}, \texttt{S10.k4Clock} and \texttt{S10.c6Clock}, for
\texttt{S10.capsule} and \texttt{S10.creg}. Lemma~\ref{lem:layer} is \texttt{S10.layer};
Lemma~\ref{lem:segments} is \texttt{S10.prefix28}, \texttt{S10.zero21}, \texttt{S10.head7},
\texttt{S10.tail11}, \texttt{S10.exit4}, \texttt{S10.rev0} and \texttt{S10.exit0};
Lemma~\ref{lem:rev} is \texttt{S10.layers}, \texttt{S10.revolution} and \texttt{S10.exit\_fires};
Theorem~\ref{thm:main} is \texttt{S10.reader\_fires\_term}, from \texttt{S10.captures\_fire},
\texttt{S10.reader\_fires\_append} and \texttt{S10.reader\_fires}, for \texttt{S10.reader},
\texttt{S10.prog}, \texttt{S10.lam}, \texttt{S10.reg} and \texttt{S10.cap}; the reader clock is
\texttt{S10.readerClock}.

\bibliographystyle{plainnat}
\bibliography{refs}
\end{document}
